\section*{ID7958: DAT-GPT Tidal Locking Criterion (Formal Statement)}
\paragraph{Metadata.}
PiID: 3420922224533 | RTEID: RTE:P30902.6813:T0.6526 | UUID: ef71ad5f-a2ce-53d4-b14d-c25f0d4267a6

\paragraph{Canonical Statement.}
\textbackslash{}[ \textbackslash{}boxed\{ \textbackslash{}left\textbackslash{}langle \textbackslash{}frac\{\textbackslash{}partial \textbackslash{}mathcal\{T\}\_\{\textbackslash{}text\{eff\}\}\}\{\textbackslash{}partial \textbackslash{}omega\} \textbackslash{}right\textbackslash{}rangle\_\{\textbackslash{}omega = p/q\textbackslash{},n\} < 0 \textbackslash{}quad\textbackslash{}text\{and\}\textbackslash{}quad \textbackslash{}frac\{\textbackslash{}Delta \textbackslash{}mathcal\{E\}\_\{\textbackslash{}text\{noise\}\}\}\{\textbackslash{}Delta t\} < \textbackslash{}left| \textbackslash{}Delta \textbackslash{}mathcal\{E\}\_\{\textbackslash{}text\{res\}\} \textbackslash{}right| \} \textbackslash{}]

\paragraph{Canonical Equation.}
\[
\boxed{ \left\langle \frac{\partial \mathcal{T}_{\text{eff}}}{\partial \omega} \right\rangle_{\omega = p/q\,n} < 0 \quad\text{and}\quad \frac{\Delta \mathcal{E}_{\text{noise}}}{\Delta t} < \left| \Delta \mathcal{E}_{\text{res}} \right| }
\]

\paragraph{Dictionary.}
DAT-GPT Tidal Locking Criterion (Formal Statement) — ⟨ frac∂ T\_eff∂ ω ⟩\_ω = p/q n < 0 and fracΔ E\_noiseΔ t < | Δ E\_res |

\paragraph{Encyclopedia.}
DAT-GPT Tidal Locking Criterion (Formal Statement) is a variational / legendre derivative, written ⟨ frac∂ T\_eff∂ ω ⟩\_ω = p/q n < 0 and fracΔ E\_noiseΔ t < | Δ E\_res |. It arises from a variational or Legendre procedure, defining conjugate quantities or stationarity conditions. It is built from ∂, ω, Δ, T, p, q. Within the Trawin Zero Point registry it serves as one element of the framework's mathematical narrative.
